\section{macOS / Linux 构建方法}

\subsection{环境准备}

\begin{enumerate}
  \item 安装 CMake（$\geq 3.24$）；
  \item 安装 arm-none-eabi 交叉工具链：
    \begin{itemize}
      \item macOS（Homebrew）：\texttt{brew install arm-none-eabi-gcc}
            或 \texttt{brew install --cask gcc-arm-embedded}；
      \item Linux：\texttt{sudo apt install gcc-arm-none-eabi}（Debian/Ubuntu）
            或使用 ARM 官方工具链包。
    \end{itemize}
\end{enumerate}

注意：Homebrew 的 \texttt{arm-none-eabi-gcc} 不带 newlib（见第
\ref{sec:restructure} 前文的修改记录），因此本工程已按「无标准库」方式配置；
若用 ARM 官方工具链（带 newlib），本工程同样可以正常编译。

\subsection{构建步骤}

在工程根目录执行：

\begin{lstlisting}[style=bash]
# 1) 删除旧构建目录（首次构建可省略）
rm -rf build

# 2) 用工具链文件配置，生成 Makefile
cmake -S . -B build -G "Unix Makefiles" \
  -DCMAKE_TOOLCHAIN_FILE=$PWD/cmake/arm-none-eabi-gcc.cmake

# 3) 编译
cmake --build build
\end{lstlisting}

\subsection{产物}

编译成功后，\texttt{build/} 下会生成：

\begin{lstlisting}[style=generic]
build/MAKE       # ELF（含调试信息）
build/MAKE.bin   # 裸二进制（烧录用）
build/MAKE.hex   # Intel HEX（烧录用）
build/MAKE.map   # 链接地址映射表
\end{lstlisting}

其中 \texttt{.bin}/\texttt{.hex} 由 \texttt{CMakeLists.txt} 里的
\texttt{POST\_BUILD} 步骤调用 \texttt{arm-none-eabi-objcopy} 自动生成。

\subsection{重新编译增量}

修改代码后只需：

\begin{lstlisting}[style=bash]
cmake --build build
\end{lstlisting}

CMake 会自动重编被改动的文件并重新链接，速度很快。

\subsection{更换生成器}

也可以用 Ninja 加速：

\begin{lstlisting}[style=bash]
cmake -S . -B build -G Ninja \
  -DCMAKE_TOOLCHAIN_FILE=$PWD/cmake/arm-none-eabi-gcc.cmake
cmake --build build
\end{lstlisting}
